\section{Introduction}
\label{sec:intro}

\noindent
This paper is the thirteenth in a series on the philosophy of intimacy and the theory of justice. Where the Prelude entered the question through a single experience, this section states it as a problem, fixes its scope, and sets out what the paper claims to establish. The question is the one the Prelude arrived at: \emc{how a promise that carries no enforcing force generates trust.} A vow, a betrothal document, a memorandum of understanding, a bare promise---an utterance that binds the one who makes it to a future conduct, but which no court will enforce, which threatens no penalty, and behind which stands no power that can compel performance---is nonetheless, in the right conditions, believed; and the believing is not credulity but a real and often well-founded expectation that the thing will be kept. The paper asks by what mechanism this expectation is produced.

The problem has a structure that makes it look, at first, impossible. Trust is directed at a future, and the future does not yet exist and cannot be verified; the symbol that asks for trust has no force, and so cannot secure the future by threatening a consequence. The thing to be explained is therefore how a symbol \emph{about what does not exist}, carrying \emph{no guarantee of consequence}, produces a real expectation in the present\,---\,something close to a generation from nothing. The paper's central and original contribution is an account of that mechanism. Its other contributions, set out below, are largely applications, extensions, and consequences of it.

\subsection*{Scope}

The paper concerns the generation of trust by non-coercive symbols, across three concentric scales: the intimate vow between two persons; the promise carried outward to a wider circle, the family in particular; and the public institutions, contract and the courts, in which the same problem appears at the scale of a society. It treats these scales as continuous: the same mechanism is argued to operate at each, and the paper moves between them rather than confining itself to one. It does not offer a general theory of trust as a psychological disposition, nor a survey of the empirical sociology of trust; it offers a mechanism for the specific case of the non-coercive promise, and traces that mechanism's reach. Where the argument requires the resources of jurisprudence, of the cognitive sciences, of psychoanalysis, of political economy, and of the Daoist and Kantian traditions, it draws on them as the prior papers in the series have drawn on them: as several irreducible epistemic positions on a single object, none of which governs the others.

\subsection*{Relation to the foundational paper and to the series}

The paper stands in a definite relation to a foundational text, to which it should be read as a sustained development of two of that text's diagnostic themes. The foundational paper~\citep{wanhong_grb} sets out the framework of which the present series is a part, and devotes nearly half its length to a diagnosis of the crisis of the symbol under modern conditions, across several dimensions. The present paper takes up two of those dimensions in particular\,---\,the crisis of authenticity and trust, and the overload of justice\,---\,and subjects them to the exhaustive mechanistic analysis that the foundational paper, whose purpose was diagnostic, deferred. Where the foundational paper established that trust has come under a crisis whose mechanism it did not specify, the present paper specifies the mechanism; where the foundational paper diagnosed an overload of the juridical and a failure of enforcement, the present paper carries that diagnosis to the claim that force has failed \emph{as a signal}, and offers the way out the foundational paper marked but did not develop. The relation to the wider series is similarly definite. The paper draws on the geometric-phase or holonomy apparatus developed earlier in the series, on its three-register theory of value, and on its account of the practitioner who cannot guarantee that his practice is good; and it observes, as the series requires, the discipline by which no single framework is permitted to absorb the others.

\subsection*{Contributions}

The paper claims to establish the following. \emph{First}, and centrally, that the expectation a non-coercive promise produces is generated by a mechanism that can be given a definite form: trust is the high precision a receiver's generative model assigns to the prediction that the promise will be kept, and this precision is supplied through three coupled channels\,---\,structural, self-committing, and wagered\,---\,which multiply rather than add, so that the consistency among them, hard to forge, is the source of well-founded trust (\xref{sec:mechanism}). \emph{Second}, a jurisprudential result: that enforcing force is not the cause of obligation but a substitute mechanism for the production of expectation, and that the internal development of legal philosophy has, over a century, moved the anchor of legal validity away from force and toward trust\,---\,a movement the law itself records, at its technical edge, in the doctrine of promissory estoppel (\xref{sec:jurisprudence}). \emph{Third}, an account of relational trust, the trust transmitted along the links of a relational network, as a phenomenon spanning all three registers, together with a new specification of how it fails under the crisis of the symbol: through a separation of the symbolic vehicle of trust from the real coupling it should carry (\xref{sec:relational}). \emph{Fourth}, a theory of the translation of sincerity across the borders of incommensurable epistemologies, with its attendant ethics (\xref{sec:translation}). \emph{Fifth}, a direction for the reform of the juridical, developed in a historical-materialist register, toward the decentralisation of enforcement (\xref{sec:judicial}). \emph{Sixth}, and as the paper's metaphysical summit, an account of self-legislation as the descent of a relational divinity, with the consequence that the promise which has no force may be, in respect of this, the more fully realised (\xref{sec:praxis}).

\subsection*{A note on method and on what is conceded}

Two commitments run through the paper. The first concerns the status of its formal and scientific apparatus. The paper draws on predictive coding for the mechanism of trust, and on the dynamics of open quantum systems for the structure of relational trust; it does so as the prior papers drew on the geometric phase, as \emph{mechanistic and structural} models and not as reductive claims about what trust ultimately is. Predictive coding furnishes a description at the computational and implementational level; it does not displace the phenomenological, the juridical, or the ethical descriptions of the same object, and it is not a meta-framework that absorbs them. The same discipline governs the use of the quantum-dynamical language in \xref{sec:relational}: a structural isomorphism is claimed, not a physical reduction. The second commitment is the series' refusal of a governing conclusion. The paper ends not in a synthesis but in plural conclusions, each framework stating what it can and cannot say, because the object\,---\,trust generated without a guarantor\,---\,is one no single framework possesses, and the form of the paper is meant to enact this as much as to assert it (\xref{sec:conclusions}).

\subsection*{Limitations and what is left open}

The paper is candid about its limits. The holonomy apparatus is employed as a structural analogy rather than developed to full formal closure here, the relevant formal programme being carried by another paper in the series; the genealogy of trust-grammars it offers (\xref{sec:translation}) is offered as adequate to its practical purpose rather than as exhaustive, and admits of finer subdivision; the depth to which the predictive-coding account is pressed, in particular toward a first-principles treatment in terms of the active reduction of expected uncertainty, is deliberately bounded, the stronger and more technical development being noted rather than carried out. The quantum-dynamical isomorphism of \xref{sec:relational} is presented as a formal analogy inviting, but not here supplying, its own development. These are marked as open rather than concealed, in keeping with the series' discipline of distinguishing what is claimed from what is conjectured.

\subsection*{Plan of the paper}

The paper proceeds as follows. It assembles, without yet explaining them, the symptoms of the contemporary condition of trust (\xref{sec:symptomatology}); and diagnoses them, drawing on the foundational paper and on the older problems of Hume and Kant, to the thesis that force is not the cause of trust (\xref{sec:diagnosis}). It then locates that thesis within the philosophy of law, tracing the migration of the anchor of legal validity from force toward trust (\xref{sec:jurisprudence}), and draws out the metaphysical consequence, that the cause of trust is inferential rather than mechanical and that this is the very opening in which freedom and the moral are possible (\xref{sec:freedom}). It develops the central mechanism (\xref{sec:mechanism}); gives its geometric reading in the language of holonomy (\xref{sec:holonomy}); states the criterion that distinguishes genuine trust from forged (\xref{sec:justice}); analyses relational trust and its failure, with its quantum-dynamical structure (\xref{sec:relational}); and develops the theory of translation by which sincerity crosses the borders of epistemology (\xref{sec:translation}). It turns then to practice: to the reform of the juridical at the scale of institutions (\xref{sec:judicial}), and to the praxis of trust at the scale of a life, rising to the account of relational divinity that is the paper's summit (\xref{sec:praxis}). It concludes in plural conclusions (\xref{sec:conclusions}), and closes in an envoi.
